\subsection{7. Environment Management: Python, R, and
Conda}\label{environment-management-python-r-and-conda}

\subsubsection{Python Experiments (most
common)}\label{python-experiments-most-common}

\begin{minted}{bash}
cd experiments/NNN_folder_name/

# Create the isolated environment
python3 -m venv .venv

# Activate (macOS/Linux)
source .venv/bin/activate

# Install dependencies
pip install -r requirements.txt

# When done working, deactivate
deactivate
\end{minted}

\begin{quote}
\small {[}IDEA{]} The \texttt{.venv/} directory is listed in
\texttt{.gitignore} --- it will NOT be committed.
That's correct. The \texttt{requirements.txt} is what
gets committed.
\end{quote}

\subsubsection{R Experiments}\label{r-experiments}

For papers using R (e.g., UBA, IRon):

\begin{minted}{bash}
cd experiments/NNN_folder_name/

# Open R and install the needed packages
Rscript -e "install.packages(c('uba', 'performanceEstimation', 'randomForest'))"

# Create an install script for reproducibility
cat > install_packages.R << 'EOF'
install.packages(c("uba", "performanceEstimation"))
EOF
\end{minted}

Document the R version in the README:

\begin{minted}{bash}
R --version | head -1
\end{minted}

\subsubsection{Conda Experiments}\label{conda-experiments}

For papers that require specific CUDA or complex binary dependencies:

\begin{minted}{bash}
conda env create -f src/original/environment.yml
conda activate experiment-name
\end{minted}
